\begin{theorem}
For each $\iota>0$, there is $D_0$ such that if $D\ge D_0$ and
$\mathcal H$ is a $3$-uniform $2$-simple hypergraph with maximum degree at
most $D$, then
\[
\chi'(\mathcal H)\le
\left(1-\frac29+\frac{2}{3^5}+\iota\right)3D.
\]
In particular, for sufficiently large $D$, every $3$-uniform $2$-simple
hypergraph with maximum degree at most $D$ has chromatic index at most
$2.3581D$.
\end{theorem}
\begin{proof}
Fix $\iota>0$.  Choose a positive real number $\alpha$ such that
$\alpha\le\iota$ and
$1-\frac29+\frac{2}{3^5}+\alpha/2<1$; this is possible because
$1-\frac29+\frac{2}{3^5}<1$.  Set
\[
\varepsilon=\alpha/2,\qquad
\theta=1-\frac29+\frac{2}{3^5}+\alpha/2 .
\]
Since $\alpha>0$, also $\varepsilon>0$, and the choice of $\alpha$ gives
$0<\theta<1$.

Apply \noderef{GeneralColoringTheorem} with $k=3$, with this
$\varepsilon$, and with the above $\theta$.  It supplies a threshold
$D_{\mathrm{GCT}}$ such that every $3$-uniform hypergraph of maximum degree at
most $D\ge D_{\mathrm{GCT}}$ is colorable with at most
$(\theta+\varepsilon)3D$ colors, provided every sub-hypergraph $H$ of maximum
degree at most $D'\ge \varepsilon D/3$ satisfies
$b_{L(H),3D'}(e)\le\theta$ for every edge $e$ of $H$.

Apply \noderef{ThreeUniformTwoSimpleBParameterEstimate} with
$\eta=\alpha/2$.  It gives a threshold $D_{\mathrm b}$ such that every
$3$-uniform $2$-simple hypergraph of maximum degree at most a natural $n\ge
D_{\mathrm b}$ satisfies
\[
b_{L(H),3n}(e)\le 1-\frac29+\frac{2}{3^5}+\alpha/2=\theta
\]
for every edge $e$.

Choose $D_0$ large enough that $D_0\ge D_{\mathrm{GCT}}$ and
$\varepsilon D/3\ge D_{\mathrm b}+1$ for every $D\ge D_0$.  Now let
$D\ge D_0$, and let $\mathcal H$ be a $3$-uniform $2$-simple hypergraph with
maximum degree at most $D$.  Consider any real $r=D'\ge \varepsilon D/3$ and
any sub-hypergraph $H\subseteq\mathcal H$ with every vertex degree at most $r$.
By \noderef{SubhypergraphInheritance}, $H$ is $3$-uniform and $2$-simple.
Set $n=\lfloor r\rfloor_+$.  The inequality $r\ge D_{\mathrm b}+1$ gives
$r\ge1$ and $n\ge D_{\mathrm b}$.  By \noderef{RealDegreeFloorBound},
$H$ has maximum degree at most $n$, so the integer b-parameter estimate
applies at $n$.  Applying \noderef{LocalBParameterRealFloorTransfer}
with $k=3$ then gives, for every edge $e$,
\[
b^{\mathbb R}_{L(H),3r}(e)\le b_{L(H),3n}(e)\le\theta.
\]

All hypotheses of \noderef{GeneralColoringTheorem} are therefore satisfied for
$\mathcal H$, with main scale $(D:\mathbb R)$.  Hence
\[
\chi'(\mathcal H)\le(\theta+\varepsilon)3D
=\left(1-\frac29+\frac{2}{3^5}+\alpha\right)3D
\le\left(1-\frac29+\frac{2}{3^5}+\iota\right)3D,
\]
which proves the first assertion.

For the decimal assertion, compute
\[
3\left(1-\frac29+\frac{2}{3^5}\right)
=\frac{3\cdot7}{9}+\frac{6}{243}
=\frac73+\frac{2}{81}
=2.358024\ldots <2.3581.
\]
Choose a positive $\iota$ so small that
$3(1-\frac29+\frac{2}{3^5}+\iota)\le 2.3581$, and let $D_0$ be the threshold
from the first assertion for this value of $\iota$.  Then for every
$D\ge D_0$ and every $3$-uniform $2$-simple $\mathcal H$ of maximum degree at
most $D$,
\[
\chi'(\mathcal H)\le
\left(1-\frac29+\frac{2}{3^5}+\iota\right)3D
\le 2.3581D .
\]
This is exactly the claimed sufficiently-large-$D$ decimal bound.
\end{proof}
